\section{线性回归与最小二乘法}\label{ch:least-squares}

\subsection{核心思路}

对于一元线性回归，预测式为 $\hat y_i=wx_i+b$。最小二乘法选择让误差平方和最小的 $w,b$：
\begin{equation}\label{eq:linear-objective}
\boxed{J(w,b)=\frac1{2m}\sum_{i=1}^{m}(wx_i+b-y_i)^2.}
\end{equation}
系数 $1/2$ 方便求导，不影响最优参数；实际报告的 MSE 则不包含这个 $1/2$。

\subsection{怎样直接求出直线}

记 $\bar x$ 为所有 $x_i$ 的平均值，$\bar y$ 为所有 $y_i$ 的平均值。对 $w,b$ 求偏导并令其为零，得到
\begin{equation}\label{eq:slope-ls}
\boxed{\hat w=
\frac{\sum_{i=1}^{m}(x_i-\bar x)(y_i-\bar y)}
     {\sum_{i=1}^{m}(x_i-\bar x)^2},
\qquad
\hat b=\bar y-\hat w\bar x.}
\end{equation}

\textbf{使用条件：}至少有两个不同的 $x_i$，否则分母为零，无法唯一确定斜率。这个公式说明拟合直线一定经过 $(\bar x,\bar y)$\cite[第3章]{james2023}。

\subsection{一个小例子}

给定 $(1,2),(2,3),(3,5)$，有 $\bar x=2,\ \bar y=10/3$。代入公式得到
\[
\hat w=\frac32,\qquad \hat b=\frac13,
\]
所以拟合直线是 $\hat y=\frac32x+\frac13$。对应的 MSE 为 $1/18$。

\subsection{需要注意什么}

最小二乘法求的是\textbf{训练数据上的最佳拟合直线}，不保证因果关系，也不保证对新数据预测得好。它不要求误差一定服从正态分布；正态假设主要用于进一步进行统计推断或似然解释。

多个特征也可以做线性回归，但会用到矩阵计算。入门阶段先掌握上述一元公式即可。